\clearpage
\begin{summarybox}

\textbf{\textcolor{CSummary}{快懂｜为什么会想到它、为什么对？}}

算 $2^x$ 与 $(\frac12)^x$ 不必每次求小数：\textbf{$x$ 每加 $1$，$a^x$ 就乘 $a$}，所以大于 $1$ 的底数让值变大，小于 $1$ 的底数让值变小；$a^0=1$ 解释了共同经过的点；$a^{-x}=(1/a)^x$ 解释了左右对称。

\medskip
\textbf{\textcolor{CSummary}{快记｜条件与最短规律}}
\begin{itemize}[leftmargin=*,itemsep=0.32em,topsep=0.3em]
  \item \textbf{先看条件：}$a>0,\ a\neq1$，指数 $x$ 可取任意实数。
  \item \textbf{正且过一：}$a^x>0$，$a^0=1$，过 $(0,1)$；值域 $(0,+\infty)$，靠近但不碰 $y=0$。
  \item \textbf{增减方向：}$a>1$ 递增；$0<a<1$ 递减。$x$ 与 $0$ 的大小决定 $a^x$ 在 $1$ 的哪边。
  \item \textbf{左右对称：}$y=a^x$ 与 $y=(1/a)^x$ 关于 $y$ 轴对称；点 $(t,y)$ 对应 $(-t,y)$。
\end{itemize}

\medskip
\textbf{\textcolor{CSummary}{快用｜考试什么时候用得最快？}}
\begin{itemize}[leftmargin=*,itemsep=0.32em,topsep=0.3em]
  \item \textbf{同底比较：}底数 $>1$，指数大者值大；底数在 $(0,1)$，指数大者值小。
  \item \textbf{解不等式：}先化同底，再比较指数；底数在 $(0,1)$ 时注意\textbf{不等号反向}。
  \item \textbf{识别图像：}先找指数为 $0$ 的位置，再看曲线的增减方向和靠近的水平线。
\end{itemize}

\textbf{\textcolor{CSummary}{最终自问：}}\textbf{为什么想到？}省去繁琐计算。\textbf{为什么成立？}因为乘法、零次方和负指数。\textbf{什么时候最快？}比较指数大小、解指数不等式、判断图像与值域。

\end{summarybox}
